\section*{Capítulo \ref{ch:b1:acetals-protection} --- Carbonilas: adições nucleofílicas, acetais e proteção}

\begin{solution}{exo:b1:acetals-protection:1}
\ce{CH3CH(OH)OCH3}: \omterm{def:b1:acetals-protection:hemiacetal}{hemiacetal}. \ce{CH3CH(OCH3)2}: \omterm{def:b1:acetals-protection:hemiacetal}{acetal}.
\ce{CH3CH(OH)2}: hidrato. \ce{CH3OCH3}: éter. \ce{(CH3)2C(OCH2CH3)2}:
\omterm{def:b1:acetals-protection:hemiacetal}{acetal} (de uma cetona).
\end{solution}

\begin{solution}{exo:b1:acetals-protection:2}
\ce{CH3CH2CHO + 2CH3OH <=> CH3CH2CH(OCH3)2 + H2O}; a propanona com o
etano-1,2-diol dá o 2,2-dimetil-1,3-dioxolano e água,
\ce{(CH3)2CO + HOCH2CH2OH <=> C5H10O2 + H2O}.
\end{solution}

\begin{solution}{exo:b1:acetals-protection:3}
O hidróxido de sódio, o brometo de metilmagnésio e o boro-hidreto de
sódio o deixam inalterado; o ácido sulfúrico aquoso diluído o hidrolisa.
\end{solution}

\begin{solution}{exo:b1:acetals-protection:4}
O OH do carbono 4 se adiciona ao carbono aldeídico 1: um anel de cinco
átomos, quatro carbonos e um oxigênio, com um OH no carbono vizinho ao
oxigênio do anel.
\end{solution}

\begin{solution}{exo:b1:acetals-protection:5}
Protonação da C=O; o metanol se adiciona ao carbono; perda de \ce{H+}:
\omterm{def:b1:acetals-protection:hemiacetal}{hemiacetal}; protonação de seu OH; perda de água até \ce{CH3CH=O+CH3}; um
segundo metanol se adiciona; perda de \ce{H+}: \ce{CH3CH(OCH3)2}. Os dois
prótons captados são devolvidos: o ácido é um \omterm{def:b1:elementary-steps:catalyst}{catalisador}.
\end{solution}

\begin{solution}{exo:b1:acetals-protection:6}
Protonação de um OCH3; perda de metanol até \ce{(CH3)2C=O+CH3}; a água se
adiciona; perda de \ce{H+}: \omterm{def:b1:acetals-protection:hemiacetal}{hemiacetal}; protonação, perda de metanol, perda
de \ce{H+}: propanona. O excesso de água mantém $Q$ abaixo de $K$ para a
hidrólise, que se completa.
\end{solution}

\begin{solution}{exo:b1:acetals-protection:7}
Uma água por \omterm{def:b1:acetals-protection:hemiacetal}{acetal}: $0.250 \times 18.0 = \qty{4.5}{g}$, cerca de
\qty{4.5}{mL}. Depois de uma hora, $3.2/18.0 = \qty{0.178}{mol}$: conversão de
\qty{71}{\percent}.
\end{solution}

\begin{solution}{exo:b1:acetals-protection:8}
Com o diol, uma molécula substitui duas: a reação não diminui o número de
moléculas (um aldeído e um diol dão um \omterm{def:b1:acetals-protection:hemiacetal}{acetal} e uma água), e o segundo OH
é mantido perto do carbono que reage, o que fecha o anel facilmente. Os
dois efeitos tornam o \omterm{def:b1:acetals-protection:hemiacetal}{acetal} cíclico mais favorável.
\end{solution}

\begin{solution}{exo:b1:acetals-protection:9}
O ácido destruiria o \omterm{def:b1:organomagnesium:organometallic}{reagente de Grignard} (transferência de próton);
traços de água e de diol também precisam ser removidos.
\end{solution}

\begin{solution}{exo:b1:acetals-protection:10}
1. Proteja a cetona da 4-clorobutan-2-ona como \omterm{def:b1:acetals-protection:hemiacetal}{acetal} cíclico
(etano-1,2-diol, \omterm{def:b1:elementary-steps:catalyst}{catalisador} ácido, Dean--Stark). 2. Magnésio em éter seco:
o \omterm{def:b1:organomagnesium:organometallic}{reagente de Grignard} do cloreto protegido, que não é mais destruído por
sua própria carbonila. 3. Adicione propanona; tratamento com cloreto de
amônio aquoso: o álcool terciário. 4. Ácido aquoso diluído: \omterm{def:b1:acetals-protection:protecting-group}{desproteção}
em 5-hidroxi-5-metil-hexan-2-ona.
\end{solution}

\begin{solution}{exo:b1:acetals-protection:11}
$Q = [\text{acetal}][\ce{H2O}]/([\text{aldeído}][\ce{CH3OH}]^2)$: um grande
excesso de metanol, elevado ao quadrado no denominador, mantém $Q < K$ até
restar pouco aldeído. Remover a água age sobre o numerador; usar o metanol
como solvente é mais simples, mas exige um produto fácil de separar.
\end{solution}

\begin{solution}{exo:b1:acetals-protection:12}
Só o OH do \omterm{def:b1:acetals-protection:hemiacetal}{hemiacetal} pode sair como água dando um cátion estabilizado
pelo oxigênio vizinho do anel (\ce{C=O+}); os outros grupos OH teriam de
sair de carbonos comuns, dando cátions não estabilizados. O metanol se
adiciona ao cátion estabilizado: um \omterm{def:b1:acetals-protection:hemiacetal}{acetal} metílico.
\end{solution}

\begin{solution}{pb:b1:acetals-protection:1}
\textbf{1.} \ce{BrMgCH2CH2CH2CHO}.
\textbf{2.} Seu carbono de caráter carbaniônico atacaria o aldeído de
outra molécula (ou dela mesma): \ce{RMgBr} se adiciona a $\mathrm{R'CHO}$, dando um \omterm{def:b1:alcohol-activation:alkoxide}{alcóxido}.
\textbf{3.} O aldeído, contra o \omterm{def:b1:organomagnesium:organometallic}{reagente de Grignard}.
\textbf{4.} Os \omterm{def:b1:acetals-protection:hemiacetal}{acetais} são estáveis frente aos \omterm{def:b1:organomagnesium:organometallic}{reagentes de Grignard} e às
bases e são removidos por ácido aquoso.
\textbf{5.} Ele se forma de modo mais completo (um diol para duas
moléculas de álcool) e é fácil de hidrolisar.
\textbf{6.} $15.09/150.9 = \qty{0.1000}{mol}$.
\textbf{7.} $1.20 \times 0.1000 \times 62.0 = \qty{7.44}{g}$.
\textbf{8.} \ce{BrCH2CH2CH2CHO + HOCH2CH2OH <=> C6H11BrO2 + H2O}
(2-(3-bromopropil)-1,3-dioxolano).
\textbf{9.} Ele ativa a carbonila e permite a perda de água na segunda
etapa; é regenerado.
\textbf{10.} O tolueno e a água destilam juntos, condensam e caem no
coletor; a água desce e fica, o tolueno transborda de volta.
\textbf{11.} $Q$ contém $[\ce{H2O}]$ no numerador; remover a água mantém
$Q < K$, e a reação avança.
\textbf{12.} A água é mais densa e não se mistura com o tolueno: ela se
acumula no fundo; só a camada superior alcança o braço de retorno.
\textbf{13.} Quando a água para de se acumular, no volume esperado.
\textbf{14.} $0.1000 \times 194.9 = \qty{19.5}{g}$.
\textbf{15.} \ce{C6H11BrO2 + Mg -> C6H11BrMgO2} (éter seco).
\textbf{16.} O reagente se adiciona a \ce{(CH3)2CO}: o \omterm{def:b1:alcohol-activation:alkoxide}{alcóxido}
\ce{(CH3)2C(O^-)CH2CH2CH2-} na cadeia protegida, com \ce{MgBr+}.
\textbf{17.} O ácido removeria o \omterm{def:b1:acetals-protection:protecting-group}{grupo protetor} cedo demais e poderia
desidratar o álcool terciário.
\textbf{18.} $0.0700 \times 174.0 = \qty{12.2}{g}$.
\textbf{19.} O \omterm{def:b1:acetals-protection:hemiacetal}{acetal} + \ce{H2O} $\to$ o aldeído + etano-1,2-diol
(\omterm{def:b1:elementary-steps:catalyst}{catalisador} ácido).
\textbf{20.} Sua perda de água exige \omterm{def:b1:predominant-reaction:strong-acid}{ácido forte} e calor; um ácido aquoso
brando e frio hidrolisa o \omterm{def:b1:acetals-protection:hemiacetal}{acetal} muito mais depressa.
\textbf{21.} O OH do carbono 5 se adiciona ao carbono aldeídico 1: um anel
de seis átomos (cinco carbonos e um oxigênio) com dois grupos metila no
carbono vizinho ao oxigênio do anel e um OH no outro carbono vizinho a ele.
\textbf{22.} $0.0700 \times 0.90 = \qty{0.0630}{mol}$, \qty{8.19}{g}.
\textbf{23.} $0.0630/0.1000 = \qty{63}{\percent}$ (\omterm{def:b1:acetals-protection:protecting-group}{proteção} considerada
completa).
\textbf{24.} $0.1000 \times 18.0 = \qty{1.80}{g}$.
\textbf{25.} $1.80/0.997 =$ \textbf{\qty{1.8}{mL} de água} no aparelho.
\end{solution}
